\FloatBarrier
\section{Discussion}
% Budget: 2 pages. Group's own work (course policy: AI may help with code, not with this component).
% Outline only: the group writes the text. Suggested structure, about half a page per point.
\todo{Discussion outline.
(1) How much of measured connectedness is a common factor? About 41\% of Diebold--Yilmaz connectedness disappears without the factor (range 38--43\% across robustness checks); one factor explains 58\% of the variance. What does this say about contagion versus common exposure (Forbes and Rigobon, 2002; Rey, 2013)? Why do the common share and DY connectedness move together, and why is that partly mechanical?
(2) Does the method change the picture? Lasso and DY almost coincide, so regularization is not what matters; NETS keeps the pattern at a lower level; FNETS measures a different object (idiosyncratic part). Which method would you use for which purpose (risk monitoring, contagion analysis, portfolio diversification)?
(3) Transmitters and receivers: South Africa, Korea, Taiwan versus Turkey, Vietnam, Argentina; how much may reflect trading hours of US-listed funds; what changes once the factor is removed.
(4) Limitations: returns instead of volatilities; ETF prices embed exchange rates and US trading hours; short windows for NETS and FNETS (positive-definite adjustment, empty idiosyncratic network in the full sample, tuning-driven drops); almost no lagged links survive in the sparse methods.}
